\documentclass{cssheet}

%--------------------------------------------------------------------------------------------------------------
% Basic meta data
%--------------------------------------------------------------------------------------------------------------

\title{Mathematische Strukturen in Okkultismus, Esoterik und Religion}
\author{Prof. Dr. Christian Spannagel}
\date{\today}
\hypersetup{%
    pdfauthor={\theauthor},%
    pdftitle={\thetitle},%
    pdfsubject={Themen Bachelorarbeit},%
    pdfkeywords={bachelor, phhd}
}

%--------------------------------------------------------------------------------------------------------------
% document
%--------------------------------------------------------------------------------------------------------------

\begin{document}

\vspace*{5mm}
\begin{center}
{\Large Themen für eine Bachelorarbeit}
\end{center}

\printtitle
\vspace*{1cm}

Mathematische Begriffe, Strukturen und Darstellungen spielen in zahlreichen religiösen, okkulten und esoterischen Traditionen eine Rolle. Beispielsweise wurden magische Quadrate, geometrische Figuren wie das Pentagramm oder bestimmte Zahlen mit besonderen, teilweise übernatürlichen Bedeutungen aufgeladen. Solche Deutungen sind historisch und kulturell bedeutsam und können zugleich interessante Anlässe für die fachmathematische und mathematikdidaktische Auseinandersetzung bieten.

Gegenstand der ausgeschriebenen Bachelorarbeiten ist die \emph{wissenschaftliche Untersuchung} solcher Themen. Es sollen die mathematischen Bezüge, historischen Entstehungsbedingungen, gesell\-schaft\-li\-chen Funktionen und didaktischen Potenziale eines esoterischen Themengebiets aus einer quellenbasierten, analytischen und kritischen Perspektive untersucht werden. Dabei sollen auch Abgrenzungen zwischen wissenschaftlicher Mathematik, religiöser Symbolik, kulturellen Deutungssystemen und pseudowissenschaftlichen Behauptungen eine wichtige Rolle spielen. Je nach Fragestellung können dabei beispielsweise fachmathematische Analysen, historische Quellenarbeit, Literaturanalysen, Schulbuch- und Lehrplananalysen oder die Entwicklung und Reflexion von Unterrichtskonzepten zum Einsatz kommen.

Mögliche Themen sind:
\begin{itemize}
\item Sternpolygone (Pentagramm, Hexagramm, Heptagramm)
\item Gematrie und Zahlenmystik in der Bibel
\item Symbole der Unendlichkeit
\item Mathematik im Volksglauben
\item Magische Quadrate
\item Heilige Geometrie
\item Arithmologie und Numerologie
\item Kryptographie in okkulten Kontexten
\item Mathematik in Horoskopen
\item Tarot
\item Bildungsplananalysen: Verortung der kritischen Diskussion von Esoterik und Potenziale für das Fach Mathematik
\item Mathematik und Mathematikdidaktik in den Schriften Rudolf Steiners
\item Analyse von Lehrplänen und Lehrmaterialien an Waldorfschulen im Fach Mathematik
\end{itemize}

Gerne dürfen eigene Schwerpunktthemen vorgeschlagen werden.

\vspace*{10mm}

\printlicense




\end{document}
